\documentclass[12pt]{article}
\usepackage[margin=2cm]{geometry}
\usepackage{color}
\usepackage{amsfonts}
\usepackage{hyperref}
\hypersetup{
    colorlinks=true,
    linkcolor=blue,
    filecolor=magenta,      
    urlcolor=blue,
    pdftitle={Overleaf Example},
    pdfpagemode=FullScreen,
    }
\pagestyle{empty}

\newcommand{\alert}[1]{{\color{red}#1}}


\begin{document}


%%%%%%%%%%%%%%%% info %%%%%%%%%%%%%%%
\fontsize{12pt}{14pt}\selectfont

\begin{flushright}
    \today
\end{flushright}
\begin{flushleft}
Editors
\\
Physical Review B
\\
American Physical Society
\end{flushleft}
\vspace{-5mm}
%%%%%%%%%%%%%%%%%%%%%%%%%%%%%%%%%%%%%

\noindent\rule{\linewidth}{1.2pt}

%%%%%%%%%%%%%%%%% main %%%%%%%%%%%%%%%
\vspace{5mm}
\fontsize{12pt}{18pt}\selectfont
\noindent Dear Editors,

\bigskip

We are pleased to resubmit our revised manuscript, “Exact Organization of Density Matrices and Entanglement Structure in the Kitaev Spin Liquid” (BS15550/Wang), for further consideration in Physical Review B. We sincerely thank you and the three Referees for their careful evaluation and constructive comments, which have helped us substantially improve both the presentation and positioning of the work.

We have carefully addressed all comments raised by the Referees. Most importantly, we have substantially revised the Introduction and Discussion to clarify the conceptual distinction between our work and previous Majorana-based calculations of the entanglement entropy. Our approach constructs the density matrix and reduced density matrices directly in the physical spin Hilbert space through their spin-string operator content. The previously known entanglement entropy therefore emerges as one consequence of the spectrum of the explicitly constructed physical RDM, rather than being the primary object of the calculation.

The physical density-matrix formulation further reveals the exact block organization and degeneracies of the RDM, their relation to gauge constraints and one-form symmetry, and a unified treatment of even- and odd-site subsystems. In response to the Referees, we have also clarified the separation between the transferable gauge-sector structure and model-dependent matter physics, expanded the physical discussion of the entanglement entropy, and addressed all technical and notational issues identified in the reports.

We believe that establishing an exact physical-spin density-matrix framework for a paradigmatic quantum spin liquid provides a useful bridge between fractionalization, emergent gauge structure, and quantum-information descriptions of quantum matter. We hope that the combination of exactness, conceptual scope, and relevance across these areas will make the revised manuscript of broad interest to the readership of Physical Review B.

Thank you for your consideration. We look forward to your decision.


\vfill
\begin{flushleft}
Sincerely,
\bigskip

Chen-Chih Wang, Yi-Ping Huang, and Sungkit Yip \\

\href{mailto:good.wangchanz@gapp.nthu.edu.tw}{good.wangchanz@gapp.nthu.edu.tw}; \href{mailto:yphuang@phys.nthu.edu.tw}{yphuang@phys.nthu.edu.tw}; \href{mailto:yip@phys.sinica.edu.tw}{yip@phys.sinica.edu.tw}
\end{flushleft}


% %%%%%%%%%%%%%%%%%%%%%%%%%%%%%%%%%%%%%%%
% \newpage

% \alert{
% \href{https://journals.aps.org/prb/authors}{PRB Publishing Guidelines}

% To ensure your papers are handled as Joint Submissions, please clearly indicate your intention in the Cover Letter of all the relevant submissions.

% The cover letter is an opportunity to explain why the manuscript is appropriate for the journal. It should include the following:

% \begin{itemize}
%     \item Context of the results
%     \item Summary of key findings
%     \item Relevant submission history at Physical Review, including another paper that is part of a joint submission
%     \item Any recommended or excluded referees
% \end{itemize}
% }




\end{document}



